\documentclass[a4paper,11pt]{article}
\usepackage[T2A]{fontenc}
\usepackage[utf8]{inputenc}
\usepackage[russian]{babel}
\usepackage{amsmath,amssymb,mathtools,cancel}
\usepackage[scaled=0.96]{helvet}
\renewcommand{\familydefault}{\sfdefault}
\usepackage{icomma}
\usepackage[left=21mm,right=21mm,top=22mm,bottom=22mm]{geometry}
\usepackage{microtype}
\usepackage{booktabs,tabularx,array}
\usepackage{enumitem}
\usepackage{tikz}
\usetikzlibrary{arrows.meta,positioning,shapes.geometric}
\usepackage{xcolor}
\usepackage{fancyhdr}
\usepackage{setspace}
\usepackage{hyperref}
\definecolor{accent}{HTML}{008080}
\definecolor{darkaccent}{HTML}{005656}
\definecolor{ink}{HTML}{303535}
\definecolor{softgray}{HTML}{666A6B}
\hypersetup{colorlinks=true,linkcolor=darkaccent,urlcolor=darkaccent}
\setlength{\parindent}{0pt}
\setlength{\parskip}{4.2pt}
\linespread{1.12}
\setlist[itemize]{leftmargin=5.5mm,topsep=3pt,itemsep=2pt,parsep=0pt}
\setlist[enumerate]{leftmargin=6mm,topsep=3pt,itemsep=6pt,parsep=0pt}
\setlength{\headheight}{23pt}
\pagestyle{fancy}
\fancyhf{}
\fancyhead[L]{\footnotesize\color{softgray}Теория вероятностей и комбинаторика}
\fancyhead[R]{\footnotesize\color{softgray}Анастасия Павлова}
\fancyfoot[C]{\footnotesize\color{softgray}\thepage}
\renewcommand{\headrulewidth}{0.35pt}
\renewcommand{\headrule}{\hbox to\headwidth{\color{accent}\leaders\hrule height \headrulewidth\hfill}}
\newcommand{\teachsection}[1]{\vspace{5pt}{\color{darkaccent}\large\bfseries #1}\par\vspace{2pt}{\color{accent}\rule{\linewidth}{0.4pt}}\par\vspace{1pt}}
\newcommand{\teachsub}[1]{\vspace{3pt}{\color{darkaccent}\normalsize\bfseries #1}\par}
\newcommand{\key}[1]{\textcolor{darkaccent}{\textbf{#1}}}
\newcommand{\solution}[1]{\par\textcolor{darkaccent}{\small\textbf{Ответ:}}\;#1}
\begin{document}
\pagenumbering{gobble}
\thispagestyle{empty}
\begin{center}
\vspace*{22mm}
{\color{darkaccent}\Large МАТЕМАТИКА}\par
\vspace{25mm}
{\fontsize{25}{29}\selectfont Чек-лист}\par
\vspace{10mm}
{\color{accent}\fontsize{20}{24}\selectfont Теория вероятностей\\[2mm] и основы комбинаторики}\par
\vspace{21mm}
{\large Подготовка к ЕГЭ по профильной математике}\par
\vfill
{\large Анастасия Павлова}\par
\vspace{13mm}
{\normalsize 19 ноября 2026 г.}\par
\vspace{21mm}
{\small Лёвин Артём Александрович\\[1mm]эксклюзивно для Анастасии Павловой}\par
\vspace{14mm}
\begin{tikzpicture}[x=1cm,y=1cm]
\draw[accent,line width=0.8pt] (0,0) .. controls (1.5,.45) and (3.0,-.45) .. (4.5,0)
.. controls (6.0,.45) and (7.5,-.45) .. (9.0,0);
\end{tikzpicture}
\end{center}
\clearpage
\pagenumbering{arabic}

\teachsection{1. Классическая вероятность: с чего начать}
Если все элементарные исходы опыта \textbf{равновозможны}, вероятность события $A$ находится по формуле
\[
\boxed{P(A)=\frac{m}{N}},
\qquad
\begin{array}{l}
N\;\text{--- число всех исходов},\\
m\;\text{--- число подходящих исходов}.
\end{array}
\]
Каждый исход учитывается \textbf{ровно один раз}. Для нескольких бросков исходом считается вся последовательность результатов, например $(2;5)$.

\teachsub{Как читать формулировку}
\begin{center}
\renewcommand{\arraystretch}{1.25}
\begin{tabularx}{\linewidth}{@{}p{0.31\linewidth}X@{}}
\toprule
\textbf{В условии} & \textbf{Что означает}\\
\midrule
Меньше $4$ & $1,2,3$; четвёрка исключена\\
Не более $4$ & $1,2,3,4$; четвёрка включена\\
Ровно один орёл & Один орёл и одна решка (в любом порядке)\\
Хотя бы один орёл & Один или несколько орлов\\
\bottomrule
\end{tabularx}
\end{center}

\teachsub{Пример 1. Один бросок правильного кубика}
При выпадении числа, \textbf{меньшего четырёх}, подходят три исхода: $1,2,3$.
\[
P=\frac{3}{6}=\frac12.
\]
Если бы требовалось число \textbf{не более четырёх}, подходили бы четыре исхода, и ответ был бы $4/6=2/3$.

\teachsection{2. Два броска кубика: метод позиций}
При двух последовательных бросках результаты записываются как \textbf{упорядоченная пара}. Всего исходов
\[
N=\underbrace{6}_{\text{первый бросок}}\cdot
\underbrace{6}_{\text{второй бросок}}=36.
\]
\textbf{Пример 2.} Первый результат равен $6$, второй произволен: $m=1\cdot 6=6$, поэтому $P=6/36=1/6$.

\textbf{Пример 3.} Оба результата одинаковы: $(1;1),\ldots,(6;6)$, то есть $m=6\cdot1=6$ и $P=1/6$.

\clearpage
\teachsection{3. Когда удобнее таблица исходов}
У двух кубиков суммы встречаются с разной частотой. В клетке таблицы записана сумма результатов двух бросков.
\begin{center}
\renewcommand{\arraystretch}{1.18}
\begin{tabular}{@{}c|cccccc@{}}
\textcolor{darkaccent}{\textbf{1-й\,$\backslash$\,2-й}} & \textbf{1} & \textbf{2} & \textbf{3} & \textbf{4} & \textbf{5} & \textbf{6}\\\hline
\textbf{1} & 2&3&4&5&6&7\\
\textbf{2} & 3&4&5&6&7&8\\
\textbf{3} & 4&5&6&7&8&9\\
\textbf{4} & 5&6&7&8&9&10\\
\textbf{5} & 6&7&8&9&10&11\\
\textbf{6} & 7&8&9&10&11&12
\end{tabular}
\end{center}

\teachsub{Пример 4. Вероятность чётной суммы}
Сумма чётна, когда оба числа \textbf{одной чётности}. Нечётных граней три, чётных тоже три:
\[
m=3\cdot3+3\cdot3=18,\qquad P=\frac{18}{36}=\frac12.
\]

\teachsub{Пример 5. Какая сумма наиболее вероятна?}
Сумму $7$ дают пары
\[(1;6),(2;5),(3;4),(4;3),(5;2),(6;1).\]
Это шесть исходов из $36$: $P(\text{сумма}=7)=1/6$. Для любой другой суммы подходящих пар меньше, поэтому \key{семь --- наиболее вероятная сумма}.

\teachsection{4. Броски симметричной монеты}
За один бросок возможны два результата: орёл (О) или решка (Р). При $k$ независимых бросках число последовательностей равно $2^k$.

\teachsub{Два броска: всё пространство исходов}
\begin{center}
\begin{tabular}{@{}cccc@{}}
\toprule
ОО & ОР & РО & РР\\\bottomrule
\end{tabular}
\end{center}
\textbf{Ровно один орёл:} ОР, РО; $P=2/4=1/2$.

\textbf{Хотя бы один орёл:} ОО, ОР, РО; $P=3/4$.

Здесь удобно применить \key{противоположное событие}:
\[
P(\text{хотя бы один орёл})=1-P(\text{ни одного орла})=1-\frac14=\frac34.
\]

\clearpage
\teachsection{5. Правило умножения и факториалы}
Если выбор происходит последовательно, количество вариантов находят \textbf{умножением числа допустимых вариантов на каждом шаге}.

\teachsub{Порядок важен}
Для $20$ различных спортсменов количество возможных порядков выступления равно
\[
20\cdot19\cdot18\cdots2\cdot1=20!.
\]
\textbf{Факториал} $n!=n(n-1)\cdots2\cdot1$, причём $0!=1$.

Если среди спортсменов $7$ россиян, вариантов с россиянином на первом месте $7\cdot19!$:
\[
P(\text{первый --- россиянин})=\frac{7\cdot19!}{20!}=\frac7{20}.
\]
Если среди $20$ участников $9$ представителей Китая, вероятность того, что \textbf{последним} выступит представитель Китая, также находится по выделенной позиции: $9\cdot19!/20!=9/20$.

\teachsub{Порядок не важен: сочетания}
Число способов выбрать $r$ человек из $N$ без учёта порядка:
\[
\boxed{C_N^{\,r}=\binom Nr=\frac{N!}{r!(N-r)!}.}
\]
\textbf{Пример 6.} Из $24$ туристов случайно выбирают троих. Какова вероятность включить заранее указанного туриста А?

Все тройки: $C_{24}^{3}$. В благоприятной тройке турист А уже выбран; двух спутников выбирают из оставшихся $23$:
\[
P=\frac{C_{23}^{2}}{C_{24}^{3}}=\frac3{24}=\boxed{\frac18}.
\]

\teachsub{Как сокращать факториалы}
\[
C_{10}^{3}=\frac{10!}{3!\,7!}
=\frac{10\cdot9\cdot8\cdot\cancel{7!}}{3\cdot2\cdot1\cdot\cancel{7!}}=120.
\]
\textbf{Важно:} $3!\,7!\neq 21!$. Нельзя перемножать числа под знаками факториала как обычные множители.

\teachsection{6. Независимые позиции: нечётное произведение}
Произведение трёх последних цифр телефонного номера нечётно, только когда \textbf{каждая} из трёх цифр нечётна. Среди $0,1,\ldots,9$ ровно пять нечётных цифр.
\[
P=\frac{5\cdot5\cdot5}{10\cdot10\cdot10}=\frac{125}{1000}=\boxed{\frac18}.
\]
Здесь цифры могут повторяться: для каждой позиции доступно десять значений.

\clearpage
\teachsection{7. Формула Бернулли: ровно $n$ успехов из $k$}
Если проводят $k$ \textbf{независимых одинаковых испытаний}, вероятность успеха в каждом равна $p$, а требуется \textbf{ровно $n$ успехов}, то
\[
\boxed{P_k(n)=C_k^{\,n}\,p^n(1-p)^{k-n}},\qquad 0\le n\le k.
\]
\begin{center}
\renewcommand{\arraystretch}{1.24}
\begin{tabularx}{\linewidth}{@{}p{0.19\linewidth}X@{}}
\toprule
\textbf{Символ} & \textbf{Что обозначает}\\\midrule
$k$ & Общее число испытаний (бросков)\\
$n$ & Требуемое число успехов\\
$p$ & Вероятность успеха в одном испытании\\
$C_k^{\,n}$ & Число способов расположить $n$ успехов среди $k$ испытаний\\\bottomrule
\end{tabularx}
\end{center}

\teachsub{Пример 7. Ровно три орла за четыре броска}
Для симметричной монеты $p=1/2$, $k=4$, $n=3$:
\[
P_4(3)=C_4^3\left(\frac12\right)^3
\left(\frac12\right)^{4-3}=4\cdot\frac1{16}=\boxed{\frac14}.
\]
Подходящих последовательностей ровно четыре: ОООР, ООРО, ОРОО, РООО. Формула совпадает с прямым перебором.

\teachsub{Пример 8. Ровно три орла за десять бросков}
\[
\begin{aligned}
P_{10}(3)&=C_{10}^3\left(\frac12\right)^3
\left(\frac12\right)^7\\
&=120\cdot\frac1{2^{10}}=\frac{120}{1024}
=\boxed{\frac{15}{128}\approx0{,}1172}.
\end{aligned}
\]
Это около $11{,}72\%$.

\teachsection{8. Частые ошибки и способы самопроверки}
\begin{itemize}
\item \key{«Ровно» и «хотя бы».} Для «хотя бы одного» включайте все случаи с одним, двумя и большим числом успехов; часто удобно дополнение до единицы.
\item \key{Порядок.} При бросках $(1;6)$ и $(6;1)$ --- разные исходы. При выборе одной и той же тройки туристов порядок перечисления участников роли не играет.
\item \key{После фиксации места.} Если один спортсмен уже выбран, оставшихся $19$ спортсменов распределяют по $19$ местам: это $19!$.
\item \key{Степени Бернулли.} Показатели $n$ и $k-n$ в сумме всегда дают $k$.
\item \key{Равновозможность.} Делить подходящие исходы на все исходы допустимо, когда исходы действительно равновозможны.
\end{itemize}

\clearpage
\begin{center}
{\color{darkaccent}\Large\bfseries Алгоритм: как выбрать способ подсчёта}
\end{center}
\vspace{11mm}
\begin{center}
\begin{tikzpicture}[
node distance=13mm and 21mm,
start/.style={ellipse,draw=accent,line width=.7pt,text width=73mm,minimum height=11mm,align=center,fill=white},
action/.style={rounded corners=3pt,draw=accent,line width=.7pt,text width=90mm,minimum height=15mm,align=center,fill=white},
decision/.style={diamond,aspect=2.4,draw=darkaccent,line width=.7pt,text width=55mm,inner sep=2pt,align=center,fill=white},
arr/.style={-{Stealth[length=2.4mm]},line width=.65pt,draw=darkaccent},
lab/.style={font=\small,fill=white,inner sep=2pt}]
\node[start] (start) {Прочитайте условие и выделите событие};
\node[decision,below=18mm of start] (ber) {Нужно ровно $n$ успехов\\из $k$ независимых испытаний?};
\node[action,below=18mm of ber] (count) {Подсчитайте все исходы\\по позициям или таблице};
\node[decision,below=16mm of count] (order) {Порядок объектов важен?};
\node[action,below left=15mm and 9mm of order,text width=49mm] (comb) {Сочетания $C_N^r$\\(выбор группы)};
\node[action,below right=15mm and 9mm of order,text width=49mm] (prod) {Правило умножения\\(последовательность)};
\node[action,below=33mm of order] (finish) {Найдите $m$ и $N$, вычислите $P=\frac{m}{N}$;\\проверьте $0\le P\le 1$};
\node[action,right=6mm of ber,text width=32mm] (bern) {Формула\\Бернулли\\$C_k^n p^n(1-p)^{k-n}$};
\draw[arr] (start) -- (ber);
\draw[arr] (ber) -- node[lab,right]{нет} (count);
\draw[arr] (ber.east) -- node[lab,above]{да} (bern.west);
\draw[arr] (count) -- (order);
\draw[arr] (order.south west) -- node[lab,above left]{нет} (comb.north);
\draw[arr] (order.south east) -- node[lab,above right]{да} (prod.north);
\draw[arr] (comb.south) |- (finish.west);
\draw[arr] (prod.south) |- (finish.east);
\end{tikzpicture}
\end{center}
\vfill
\begin{center}
\small\color{softgray}Выбор метода зависит от формулировки события, числа испытаний и того, важен ли порядок.
\end{center}

\clearpage
\teachsection{Тренировочный блок}
\textcolor{softgray}{Путь А \;\textbullet\; 10 заданий. Решайте по порядку: от базовых подсчётов к сочетаниям и условию на две позиции.}
\begin{enumerate}[label=\textcolor{darkaccent}{\textbf{\arabic*.}}]
\item Правильный игральный кубик бросают один раз. Найдите вероятность выпадения числа \textbf{не более четырёх}.
\item Кубик бросают дважды. Найдите вероятность того, что в \textbf{первом} броске выпадет пятёрка.
\item Кубик бросают дважды. Найдите вероятность \textbf{нечётной суммы} выпавших очков.
\item Кубик бросают дважды. Найдите вероятность того, что сумма очков равна \textbf{восьми}.
\item Симметричную монету бросают трижды. Найдите вероятность того, что выпадет \textbf{ровно два орла}.
\item Симметричную монету бросают четыре раза. Найдите вероятность того, что выпадет \textbf{хотя бы одна решка}.
\item Симметричную монету бросают восемь раз. По формуле Бернулли найдите вероятность выпадения \textbf{ровно трёх орлов}.
\item Последние четыре цифры номера выбираются независимо и равновероятно из цифр от $0$ до $9$. Найдите вероятность того, что \textbf{их произведение нечётно}.
\item Из $18$ участников случайным образом выбирают команду из трёх человек. Найдите вероятность включения в неё \textbf{заранее указанного участника}.
\item В соревновании участвуют $20$ спортсменов: $7$ из России, $9$ из Китая и $4$ из других стран. Порядок выступлений случаен. Найдите вероятность того, что \textbf{первым выступит россиянин, а последним --- спортсмен из Китая}.
\end{enumerate}
\vspace{5mm}
\textcolor{darkaccent}{\textbf{Перед проверкой ответов}}\quad Укажите для каждой задачи число исходов, число подходящих случаев и почему они равновозможны. В задачах 7 и 9 покажите работу с факториалами.

\clearpage
\teachsection{Ответы к тренировочному блоку}
\begin{center}
\renewcommand{\arraystretch}{1.50}
\begin{tabularx}{\linewidth}{@{}p{16mm}p{46mm}X@{}}
\toprule
\textbf{№} & \textbf{Ответ} & \textbf{Контрольный подсчёт}\\\midrule
1 & $\dfrac23$ & $4/6$\\
2 & $\dfrac16$ & $(1\cdot6)/36$\\
3 & $\dfrac12$ & $18/36$\\
4 & $\dfrac5{36}$ & $(2;6),(3;5),(4;4),(5;3),(6;2)$\\
5 & $\dfrac38$ & $C_3^2/2^3$\\
6 & $\dfrac{15}{16}$ & $1-1/2^4$\\
7 & $\dfrac7{32}$ & $C_8^3/2^8=56/256$\\
8 & $\dfrac1{16}$ & $5^4/10^4$\\
9 & $\dfrac16$ & $C_{17}^2/C_{18}^3$\\
10 & $\dfrac{63}{380}$ & $(7\cdot9)/(20\cdot19)$\\
\bottomrule
\end{tabularx}
\end{center}

\vspace{9mm}
\teachsection{Итоговая памятка}
\begin{center}
\renewcommand{\arraystretch}{1.35}
\begin{tabularx}{\linewidth}{@{}p{0.38\linewidth}X@{}}
\toprule
\textbf{Сигнал в задаче} & \textbf{Полезный приём}\\\midrule
Один бросок / один случайный выбор & Перечислите равновозможные исходы\\
Несколько независимых позиций & Правило умножения\\
Два кубика; сумма значений & Таблица $6\times6$ или подсчёт пар\\
Хотя бы один успех & Противоположное событие\\
Ровно $n$ успехов из $k$ & Формула Бернулли\\
Группа без учёта порядка & Сочетания $C_N^r$\\
Первый или последний в очереди & Зафиксируйте нужное место\\
\bottomrule
\end{tabularx}
\end{center}
\vfill
\begin{center}
{\footnotesize\color{softgray}Лёвин Артём Александрович \;\textbullet\; эксклюзивно для Анастасии Павловой}
\end{center}
\end{document}
